% BAN TIENG VIET cua 04_anchor.tex (2026-10-08). So lieu giu nguyen.
\section{Chứng thực neo theo thiết bị}
\label{sec:anchor}

\subsection{Bản ghi của gateway}

Các mệnh đề~\ref{prop:blind}--\ref{prop:coincide} cho thấy tổng lượng và độ
sâu là chưa đủ, và thứ tự thời gian chính xác là thứ đối thủ thao túng, nên
phòng thủ buộc phải quan sát dấu thời gian --- và bên duy nhất có thể bảo chứng
cho một dấu thời gian là thiết bị đã thực thi nó. Giả sử gateway giữ một bản
ghi lệnh có chữ ký
\begin{equation}
\mathcal{G}=\{(t_j,u_j,\sigma_j)\},
\label{eq:gateway}
\end{equation}
trong đó $\sigma_j$ là chữ ký trên $(t_j,u_j)$ do bộ điều khiển tạo ra tại
thời điểm đóng mở máy bơm. Phép thử neo đối chiếu nhật ký khai báo với bản ghi
đó:
\begin{equation}
\begin{aligned}
\Delta_{\mathrm{time}}&=\{t_i\}\triangle\{t_j\},\\
\Delta_{\mathrm{depth}}&=\max_{t}\bigl|u(t)-u_{\mathcal{G}}(t)\bigr|,
\end{aligned}
\label{eq:deltas}
\end{equation}
trong đó $\triangle$ là hiệu đối xứng và $\Delta_{\mathrm{depth}}$ lấy trên các
giờ có mặt ở cả hai bản ghi. Phán quyết khi đó là
\begin{equation}
A(L,\mathcal{G})=
\begin{cases}
\text{giả mạo}, & \Delta_{\mathrm{time}}\neq\emptyset
\ \vee\ \Delta_{\mathrm{depth}}>\epsilon,\\
\text{khớp}, & \text{trường hợp còn lại},
\end{cases}
\label{eq:anchor}
\end{equation}
với $\epsilon=10^{-6}$~mm, tức khớp chính xác tới mức nhiễu dấu phẩy động.
Việc gọi tên riêng hai loại sai lệch là quan trọng: dịch thời điểm chỉ làm thay
đổi $\Delta_{\mathrm{time}}$, thổi phồng độ sâu chỉ làm thay đổi
$\Delta_{\mathrm{depth}}$, còn bỏ sót và thêm mới làm thay đổi cả hai, nên một
phép đối chiếu duy nhất bắt được mọi lớp gian lận.

\begin{proposition}[Tính đầy đủ của neo thiết bị]
\label{prop:complete}
Nếu nhật ký khai báo $L$ khác bản ghi gateway $\mathcal{G}$ về đa tập hợp các
giờ hoặc về bất kỳ độ sâu nào tại một giờ chung, thì $A(L,\mathcal{G})$ trả về
\text{giả mạo}. Nếu $L$ và $\mathcal{G}$ khớp cả hai, thì $L$ là một bản sao
trung thực và không có sự giả mạo nào đối với ba đại lượng trong
\eqref{eq:three}.
\end{proposition}

\begin{proof}
Xem Phụ lục~S12: mệnh đề thứ nhất là định nghĩa của \eqref{eq:anchor}, còn
mệnh đề thứ hai suy ra được vì cả ba đại lượng có thể làm giả đều là hàm của
tập giờ và các độ sâu theo giờ.
\end{proof}

\subsection{Vì sao điều này không đưa việc kết tội oan trở lại}

Mệnh đề~\ref{prop:complete} mạnh tới mức trông nguy hiểm: mọi sai lệch đều bị
coi là giả mạo, kể cả việc quên ghi và làm tròn mà Mục~\ref{sec:guarantee} đã
chỉ ra là rất phổ biến. Cách hoá giải là ở chỗ neo thiết bị đối chiếu một bản
sao với \emph{nguồn}, chứ không đối chiếu một nhật ký với một \emph{mô hình}:
một nông dân trung thực quên một mục ghi chỉ đơn giản được bổ sung từ
$\mathcal{G}$. Khi đó, kết tội oan đòi hỏi gateway phải mâu thuẫn với chính nó
--- một vấn đề toàn vẹn thiết bị --- đó là lý do Bảng~\ref{tab:falseacc} giữ ở
mức $0/60$ trong khi các tầng hành vi giữ ở mức $26{,}7\%$ và $5{,}0\%$
(Phụ lục~S18).

\subsection{Chi phí và phạm vi của tuyên bố}

Neo thiết bị đòi hỏi gateway của công trình đồng hành~\cite{twingate2026}, vốn
đóng mở máy bơm và đã ghi một bản ghi lệnh có chữ ký kèm dấu thời gian như một
phần chức năng an toàn của nó; chi phí biên của việc chứng thực là lưu trữ và
truyền một bản ghi đằng nào cũng tồn tại (Bảng~S5 của Phụ lục). Dù vậy, phạm
vi hẹp hơn tiền đề ban đầu của chúng tôi: một nhật ký viết tay, không có thiết
bị nào trong vòng lặp, tự nó \emph{không} phải bằng chứng được chấp nhận, vì
Mục~\ref{sec:results-attack} đo được một đối thủ duy lý thổi phồng tín chỉ thêm
$+35\%$ trong khi lọt qua mọi phép thử dựa trên vật lý và tổng lượng ($6/6$
nhật ký bị sửa). MRV với chi phí gần như bằng không chỉ khả thi ở nơi có một bộ
điều khiển đứng giữa nông dân và máy bơm; ở nơi tưới thủ công, lớp này không áp
dụng được.
